%======================================================================
\section{Calendar: order the ledgers}
\label{sec:calendar}

One dimension has been silent: time.  \Cref{prop:noarb} silences
butterflies within a slice, but a surface is many slices, and fitting
them independently leaves a second static arbitrage on the table.  The
requirement is a comparison of normalized calls at equal
\emph{log-moneyness},
\begin{equation}
c_1(k)\;\le\;c_2(k)
\qquad\text{for all }k,\qquad \tau_1<\tau_2 ,
\label{eq:calreq}
\end{equation}
and it must be derived, not asserted, because at equal $k$ the two
options are struck at \emph{different dollar strikes}
($K_i=F_ie^{k}$), so \eqref{eq:calreq} is not the plain listed calendar
spread it may evoke.  The derivation is three lines once the standing
assumptions are made to work in the open.

\begin{lemma}[Calendar ordering at equal log-moneyness]
\label{lem:caljensen}
Assume deterministic rates and carry, so that all forward measures
coincide with one risk-neutral measure $\Prob$, today's forward curve
$t\mapsto F_{0,t}$ is deterministic, and the deflated spot
$M_t=S_t/F_{0,t}$ is a positive $\Prob$-martingale with $M_0=1$.  Then
$M_{T_i}=S_{T_i}/F_i=Y_i$ is exactly the normalized return of slice
$i$, and for $T_1<T_2$ and every $k$,
\[
c_2(k)=\E\big[(M_{T_2}-e^{k})^{+}\big]
=\E\Big[\E\big[(M_{T_2}-e^{k})^{+}\,\big|\,\mathcal F_{T_1}\big]\Big]
\;\ge\;\E\Big[\big(\E[M_{T_2}\mid\mathcal F_{T_1}]-e^{k}\big)^{+}\Big]
=\E\big[(M_{T_1}-e^{k})^{+}\big]
=c_1(k),
\]
by the tower property, conditional Jensen applied to the convex payoff,
and the martingale property.
\end{lemma}

The deterministic-carry assumption is doing all the visible work: it is
what makes one deflator serve both expiries, hence what makes the
\emph{normalized} calls comparable at equal $k$ at all.  The
monetizing trade behind a violation is correspondingly
\emph{forward-scaled} in strike --- short $1/(D_1F_1)$ calls struck at
$K_1=F_1e^{k}$ against long $1/(D_2F_2)$ calls struck at
$K_2=F_2e^{k}$ --- not the equal-dollar-strike listed spread; the
classical statement of the static conditions in these units is in
\citet{CarrMadan2005}.  In the language of probability,
\eqref{eq:calreq} together with equal means says the laws of $Y$
increase in \emph{convex order}: the longer-dated law is worth at least
as much against every convex payoff.  Checking a continuum of call
inequalities sounds like work.  It is not, because the call curve and
the ledger are conjugates of each other, and pointwise domination
passes through conjugacy in both directions.

\subsection{Conjugate duality, and the order theorem}

\begin{lemma}[Ledger--call conjugacy]
\label{lem:conjugacy}
For any admissible slice, with $u=\Lambda(z)$ throughout,
\begin{equation}
c(k)=\max_{z\in[-\infty,\infty]}\Big[G(z)-e^{k}\big(1-\Lambda(z)\big)\Big],
\qquad
G(z)=\min_{k\in[-\infty,\infty]}\Big[c(k)+e^{k}\big(1-\Lambda(z)\big)\Big],
\label{eq:conjugacy}
\end{equation}
the maximum attained at $z=z_k$ and the minimum at $k=x(z)$.
\end{lemma}

\begin{proof}
For any $z$,
\[
G(z)-e^{k}\big(1-\Lambda(z)\big)
=\int_z^{\infty}\big(e^{x(t)}-e^{k}\big)\,\rho(t)\,\dd t ,
\]
whose integrand is negative exactly for $t<z_k$ and positive exactly for
$t>z_k$ (monotonicity of $x$).  Starting the integral before $z_k$ adds
negative mass; starting it after $z_k$ forgoes positive mass; so the
expression is maximal at $z=z_k$, where it equals $c(k)$ by
\eqref{eq:call}.  (The endpoints $z=\pm\infty$ give $1-e^{k}\le c(k)$ and
$0\le c(k)$, consistent.)  For the second identity, fix $z$.  By the
first identity, for every $k$,
\[
c(k)+e^{k}\big(1-\Lambda(z)\big)
\;\ge\;\Big[G(z)-e^{k}\big(1-\Lambda(z)\big)\Big]
+e^{k}\big(1-\Lambda(z)\big)=G(z),
\]
with equality at $k=x(z)$, where the maximizing $z'$ of the first
identity is $z$ itself.  As $k\to-\infty$ the expression tends to
$\E[Y]=1=G(-\infty)\ge G(z)$ and as $k\to+\infty$ it diverges, so the
minimum over the closed line is attained at the finite root.
\end{proof}

The first identity of \eqref{eq:conjugacy} is the pricing formula
\eqref{eq:call} rediscovered as a maximization; the pair says the call
curve and the ledger are (essentially) Legendre transforms of one
another, with the strike root as the point of tangency.

\begin{theorem}[Calendar order is ledger order]
\label{thm:ledgerorder}
Let two admissible slices have call curves $c_1$, $c_2$ and ledgers
$G_1$, $G_2$ (both mean one).  Then
\[
c_1(k)\le c_2(k)\ \ \text{for all }k
\qquad\Longleftrightarrow\qquad
G_1(z)\le G_2(z)\ \ \text{for all }z .
\]
\end{theorem}

\begin{proof}
($\Leftarrow$) For every $k$, using the max-representation and
$G_1\le G_2$ pointwise,
\[
c_1(k)=\max_z\Big[G_1(z)-e^{k}\big(1-\Lambda(z)\big)\Big]
\le\max_z\Big[G_2(z)-e^{k}\big(1-\Lambda(z)\big)\Big]=c_2(k) .
\]
($\Rightarrow$) For every $z$, using the min-representation and
$c_1\le c_2$ pointwise,
\[
G_1(z)=\min_k\Big[c_1(k)+e^{k}\big(1-\Lambda(z)\big)\Big]
\le\min_k\Big[c_2(k)+e^{k}\big(1-\Lambda(z)\big)\Big]=G_2(z) .
\qedhere
\]
\end{proof}

A continuum of option inequalities has become the pointwise ordering of
two curves that are already on the grid --- the same array that priced
the calls and carried the Jacobian, doing its fifth job.  Note where the
mean-one normalization enters: both ledgers start at $G(-\infty)=1$, so
the ordering is of two curves pinned together at one end, and a
legitimate term structure clears the bar by a thin margin everywhere ---
which is why calendar violations are easy to manufacture and hard to see
by eye on smile plots.

\begin{remark}
Pointwise call domination with equal means is the classical
characterization of convex order (calls plus affine payoffs generate the
convex functions), and by the theorems of Strassen and Kellerer
\citep{Strassen1965,Kellerer1972} a family of laws increasing in convex
order is exactly one that some martingale can run through.
\Cref{thm:ledgerorder} is thus the model's native test for ``these two
slices could belong to one arbitrage-free market''.
\end{remark}

\subsection{Enforcement is a control, not a theorem}
\label{sec:calcontrol}

Within one slice, validity was structural.  \emph{Across} slices it is
not: the family is closed under nothing that relates two different
expiries, so calendar order must be imposed by the objective, and the
paper's obligation is to say exactly what is and is not guaranteed.  The
production scheme adds, when fitting a far slice against an
already-fitted nearer one, hinge rows
\begin{equation}
\sqrt{w_{\mathrm{cal}}}\;\zeta_m\,
\big(c_{\mathrm{near}}(k_m)-c_{\mathrm{far}}(k_m)\big)^{+},
\label{eq:calrows}
\end{equation}
at fixed strikes $k_m$ --- with three deliberate scope decisions.

\emph{Price space, not ledger space.}  \Cref{thm:ledgerorder} makes the
ledger comparison the natural \emph{measurement}, and it remains the
diagnostic of record: rank space needs no volatility inversion and keeps
meaning in the far wings, exactly where inversion goes blind.  But as a
\emph{constraint} the ledger is a trap: $G(z)$ integrates the entire
upper tail, so a ledger floor at any single rank drags both slices'
extrapolated wings --- curve regions no one quoted --- into the
constraint.  On an acute short-dated pair, a full-grid ledger floor
dragged the far expiry's \emph{worst quote error} from
\MacPhantomDragFromBp\ to \MacPhantomDragToBp\ vol bp, a
two-order-of-magnitude phantom cost that vanished when the same
ordering was read at fixed strikes.  The constraint rows
\eqref{eq:calrows} therefore compare prices, which are local in
strike.

\emph{Support confinement, with a smooth taper.}  The strikes $k_m$ run
over the \emph{intersection of the two slices' retained quote spans}
only, and the weights $\zeta_m$ taper to zero in a $\cos^2$ ramp
(continuously differentiable, so the solver sees no kink) over the outer
$15\%$ of that span.  Outside the common quoted support the constraint
is silent by design: enforcing order between two extrapolations would
let one model prior push another around in a region where the market
said nothing.  The count of constraint strikes scales with basis order
($\max(49,\,4N+1)$ nodes), so a high-order slice cannot slip a violation
between nodes of a fixed coarse grid.

\emph{Finite weight, honest reporting.}  The rows carry a large but
finite weight; a hard-pressed fit can retain a small violation rather
than destroy its quote fit.  Residual violations are measured and
reported in economic units --- worst-violation strike, normalized-price
gap, and the gap as a fraction of the local bid--ask width --- so that
``violation of $10^{-6}$'' resolves into ``a fortieth of the spread:
unharvestable'' or into a number that stops publication.

The guarantee, stated plainly: \emph{within the common quoted support,
calendar order holds up to a reported, economically sized tolerance;
outside it, calendar order is not enforced, and the ledger diagnostic
reports whatever the priors did}.  This is a control, not a theorem, and
the distinction is visible in live data.

\paragraph{A live illustration.}
In the frozen snapshot (\cref{app:data}), the two shortest SPY expiries
--- two and four \emph{calendar} days from the reference date --- are
fitted cleanly, each valid on its own, calendar rows active on their
common support.  Yet the pair trips a \emph{vol-space} calendar audit
run at equal log-moneyness over the full display grid, both legs' total
variances read as volatilities on the far node's clock: a crossing of
\MacCalGapVolBp\ vol bp, argmax at $k=\MacCalGapArgmaxK$.  The two
slices' common quote span is $[\MacCalSpanLo,\,\MacCalSpanHi]$: the
flagged strike sits dozens of span-widths to the right of anything
quoted, where the two call prices are \MacCalendarCallShort\ (two-day
leg) and \MacCalendarCallLong\ (four-day leg) of the forward.  A
price-space audit over the same grid finds a worst violation of
\MacCalPriceGridWorst\ of forward --- round-off, not economics.  What
the vol-space audit detected is real --- the two \emph{extrapolations}
cross, in a region where implied variance is astronomically sensitive
to price --- and harvesting it would mean trading options dozens of
span-widths out of the money at those prices: there is no arbitrage
there, only two priors disagreeing about a wilderness.
\Cref{fig:calendar} draws the pair, the common span, and the crossing;
the paper's position is that a fitter should enforce where the market
speaks, report honestly where it does not, and resist the temptation to
make the wilderness look tidy by constraining it.

\begin{figure}[htbp]
\centering
\paperfig{\textwidth}{fig_calendar.pdf}
\caption{Calendar order as control, in two panels.  Panel A: SPY
total-variance term structure at several fixed log-moneyness values
from the frozen snapshot --- monotone through the quoted region, as the
constraint rows \eqref{eq:calrows} require.  Panel B: the confinement
illustration on the two- and four-calendar-day pair, both curves drawn
as $\sqrt{w(k)/\tau_{\mathrm{far}}}$ --- total variances annualized on
the common far-node clock, so the vertical axis reads in volatility
units and the audit gap is the visible gap between the curves.  The
common quote span is shaded, and within it the curves are ordered; each
solid curve is truncated where its node's call price drops below
$10^{-13}$ of forward and continued as a light dotted flat --- beyond
those points the stored smile is Black-inversion noise, which is itself
the argument for support-confined enforcement.  The audit crossing at
$k=\MacCalGapArgmaxK$ (gap \MacCalGapVolBp\ vol bp, call prices
\MacCalendarCallShort\ and \MacCalendarCallLong\ of forward, annotated)
lies in that censored region: a disagreement between priors in an
economically empty region, reported, not enforced away.}
\label{fig:calendar}
\end{figure}
